\subsection*{Spelling from additional Chinese image sequences}
We screened 271,574 original training records and retained one distinct, fully renderable human answer of 6--96 glyphs per accepted record. Exclusions covered the preceding fit and held-out record identifiers, normalized held-out questions, exact held-out answers and 16-character matches to the held-out questions, references or retrieved paragraphs. Screening and rendering took place offline. The resulting sequences contained 7,725 distinct bitmaps. The ordered and shuffled indexes had identical lengths and unigram frequencies. Corpus storage and full source locators remain local under the DuReader research terms.

At each writing decision, the model began with its neural distribution over training-answer bitmaps. For suffixes of one to four emitted glyphs, it mixed in the empirical next-image distribution with weight $\min(0.8,n/(n+8+d))$, where $n$ counts eligible successors and $d$ counts distinct eligible shapes. An unseen suffix contributed zero weight. The writer chose the most probable complete bitmap. This restricted study tests familiar-shape spelling. The autonomous comparison reproduced the complete preceding finite-writer pixel archive, then scored the ordered, shuffled and rotated-question conditions. A separate calculation recounted 2,284 complete streams and the 1,525 supplied-history positions from stored physical outputs and decisions. Index construction took 60.1 seconds of wall time. Scoring the neural and empirical combinations took 11.4 seconds on one RTX 4090 D and used 56.8\,MiB of peak PyTorch allocation.
